\documentclass{article}
\usepackage{amsmath}
\usepackage{amssymb}
\usepackage{geometry}
\geometry{margin=1in}

\title{Modelo de Produccion con Inteligencia Artificial}
\author{Sistema Multi-Agente de Investigacion}
\date{Abril 2026}

\begin{document}

\maketitle

\section{Funcion de Produccion Modificada}

Proponemos una modificacion de la funcion de produccion Cobb-Douglas para incorporar efectos de la IA en la industria del software:

\begin{equation}
Y = A \cdot K^{\alpha} \cdot L^{\beta} \cdot (1 + A_i)^{\gamma} \cdot (1 - V_c)^{\delta}
\end{equation}

Donde:
\begin{itemize}
    \item $Y$ = Output de software (lineas de codigo productivas, features entregadas, valor generado)
    \item $A$ = Factor de productividad total (TFP)
    \item $K$ = Capital (infraestructura, licencias de software, compute)
    \item $L$ = Trabajo (horas de desarrolladores, cantidad de developers)
    \item $A_i$ = Aumento de productividad por IA (rango tipico: 0.2 - 4.0)
    \item $V_c$ = Vulnerabilidad por vibe coding (rango: 0.0 - 0.8)
    \item $\alpha, \beta, \gamma, \delta$ = elasticidades de produccion
\end{itemize}

\section{Interpretacion de las Variables de IA}

\subsection{Variable $A_i$ - Amplificacion por IA}

$A_i$ representa el multiplicador de productividad que la IA proporciona:

\begin{equation}
A_i = f(M, T, E, Q)
\end{equation}

Donde:
\begin{itemize}
    \item $M$ = Capacidad del modelo de IA utilizado
    \item $T$ = Training/quality assurance invertido
    \item $E$ = Experiencia del developer con herramientas de IA
    \item $Q$ = Calidad de los prompts y arquitectura
\end{itemize}

\subsection{Variable $V_c$ - Vulnerabilidad por Vibe Coding}

$V_c$ captura la friccion por deuda tecnica generada por codigo de baja calidad:

\begin{equation}
V_c = \frac{D_{acumulada}}{D_{critica}} \cdot \frac{B_{produccion}}{B_{tolerable}}
\end{equation}

Donde:
\begin{itemize}
    \item $D_{acumulada}$ = Deuda tecnica acumulada (en anos-hombre)
    \item $D_{critica}$ = Umbral critico de deuda tecnica
    \item $B_{produccion}$ = Bugs en produccion por 1000 LOC
    \item $B_{tolerable}$ = Umbral tolerable de bugs
\end{itemize}

\section{Elasticidades Estimadas}

Basado en datos empiricos 2024-2026:

\begin{table}[h]
\centering
\begin{tabular}{|c|c|c|}
\hline
\textbf{Parametro} & \textbf{Valor} & \textbf{Interpretacion} \\
\hline
$\alpha$ (capital) & 0.30 & Elasticidad del capital \\
$\beta$ (labor) & 0.50 & Elasticidad del trabajo \\
$\gamma$ (IA) & 0.40 & Elasticidad de la amplificacion IA \\
$\delta$ (deuda) & 0.60 & Elasticidad de la deuda tecnica \\
\hline
\end{tabular}
\end{table}

Nota: $\delta > \gamma$ indica que la deuda tecnica tiene mayor impacto negativo que el beneficio positivo de la IA.

\section{Derivacion de la Tasa de Cambio}

Tomando logaritmos y derivando respecto al tiempo:

\begin{equation}
\frac{d \ln Y}{dt} = \frac{d \ln A}{dt} + \alpha \frac{d \ln K}{dt} + \beta \frac{d \ln L}{dt} + \gamma \frac{d \ln (1 + A_i)}{dt} - \delta \frac{d \ln (1 - V_c)}{dt}
\end{equation}

Simplificando:

\begin{equation}
g_Y = g_A + \alpha g_K + \beta g_L + \gamma \frac{\dot{A_i}}{1 + A_i} - \delta \frac{\dot{V_c}}{1 - V_c}
\end{equation}

Donde $g_X$ representa la tasa de crecimiento de la variable $X$.

\section{Condicion de Sostenibilidad}

Para que el crecimiento sea sostenible a largo plazo:

\begin{equation}
\gamma \frac{\dot{A_i}}{1 + A_i} > \delta \frac{\dot{V_c}}{1 - V_c}
\end{equation}

Esto implica que la mejora en la calidad de uso de la IA debe superar el crecimiento de la deuda tecnica, ponderado por sus respectivas elasticidades.

\section{Extension: Modelo con Segmentacion de Usuarios}

Para capturar los diferentes segmentos de usuarios (Casual, Vibe Coders, Power Users), extendemos el modelo:

\begin{equation}
Y_{total} = \sum_{j \in \{C, V, P\}} w_j \cdot Y_j
\end{equation}

Donde:
\begin{itemize}
    \item $C$ = Casual Users (45\%)
    \item $V$ = Vibe Coders (35\%)
    \item $P$ = Power Users (20\%)
    \item $w_j$ = pesos de participacion en el mercado
\end{itemize}

Cada segmento tiene sus propios parametros:

\begin{equation}
Y_j = A \cdot K^{\alpha} \cdot L^{\beta} \cdot (1 + A_{i,j})^{\gamma_j} \cdot (1 - V_{c,j})^{\delta_j}
\end{equation}

\begin{table}[h]
\centering
\begin{tabular}{|c|c|c|c|}
\hline
\textbf{Segmento} & $\gamma_j$ & $\delta_j$ & $A_{i,j}$ \\
\hline
Casual & 0.15 & 0.20 & 0.3 \\
Vibe Coders & 0.35 & 0.80 & 1.2 \\
Power Users & 0.60 & 0.15 & 3.0 \\
\hline
\end{tabular}
\end{table}

\section{Conclusion}

Este modelo captura la dualidad de la IA en el desarrollo de software: potencial amplificador ($A_i$) versus riesgo de deuda tecnica ($V_c$). Los datos empiricos sugieren que $\delta > \gamma$, indicando que sin gobernanza adecuada, los costos de la deuda tecnica pueden superar los beneficios de productividad.

\end{document}
